% !PRINTED 180
\weberpage{180}

%% [running head, not body text -- commented out 2026-08-07] Page furniture (running
%% head, folio, signature) is excluded from this edition by standing policy.
%% \noindent 180 \hfill \textit{Sechster Abschnitt.} \hfill \S.\,46.

%<<CONT
changed. $s^{-1}\tau s$ is therefore also contained in $T$, and $T$ is consequently a normal divisor of $S$.

We denote by $\mu, \nu$ the degrees of $S$ and $T$ and set
\[
\mu = e\nu;
\]
then $e$ is the index of $T$ and at the same time the degree of the group $\Sigma$.

Every absolute invariant of the group $T$, i.\,e. every function of $x$ that admits the substitutions of $T$, can be expressed rationally in $\Omega$ in terms of the $X$.

This follows by the method of inference already often applied: if we denote by $\varrho$ a function of the $X$ that takes $e$ distinct values $\varrho, \varrho_1, \ldots \varrho_{e-1}$, and set
\weq{2}{\Phi(t) = (t-\varrho)(t-\varrho_1)\ldots(t-\varrho_{e-1})}
then $\Phi(t)$ is a rational function of the variable $t$ in $\Omega$.

An absolute invariant $r$ of $T$ takes at most $e$ distinct values, which correspond to the values $\varrho, \varrho_1, \ldots \varrho_{e-1}$, namely $r, r_1, \ldots r_{e-1}$, where in certain circumstances equal values may also occur. The function of $t$:
\[
\frac{\Phi(t)\,r}{t-\varrho} + \frac{\Phi(t)\,r_1}{t-\varrho} + \cdots + \frac{\Phi(t)\,r_{e-1}}{t-\varrho_{e-1}} = \Psi(t) % [sic] second denominator printed t-\varrho, not t-\varrho_1
\]
is then contained in $\Omega$, and for $t=\varrho$ we obtain:
\weq{3}{r = \frac{\Psi(\varrho)}{\Phi'(\varrho)},}
whereby, since $\Phi'(\varrho)$ is different from zero, $r$ is rationally expressed by $\varrho$. Thus $r$ is contained in the field $\Omega(X_1, X_2, \ldots X_m)$.

The invariants of the group $T$ are at the same time invariants of the group $S$; by the solution of the form problem for the group $\Sigma$ the $X_1, \ldots X_m$ are then known, and hence also the invariants of the group $T$, and the form problem of the group $S$ is thus reduced to the successive solution of the two form problems of the groups $\Sigma$ and $T$, which are of lower degree than the form problem for $S$. A form problem which is decomposable in this way into two form problems of lower degree we may call an \gesperrt{imprimitive form problem}. Since for such a reduction it is necessary that $T$ be a normal divisor of $S$ different from $S$ and from the unit, the corresponding form problem, when $S$ is a simple group, is always primitive.
